\silentchapter{8 Stolen}{stolen}

\begin{overgang}
Om lag 2\,300 europeiske stolar, laga mellom 1280 og 2024, utgjer den empiriske ryggraden i denne boka. Dei kjem frå Nasjonalmuseet i Oslo og Victoria and Albert Museum i London. Kvar stol er målt, katalogisert og plassert i formrommet. Det som følgjer, er historia om kva stolane viser oss.
\end{overgang}

\section*{Formrommet er ikkje uniformt busett}

Dersom alle seleksjonstrykk var like sterke og alle materiale like tilgjengelege over heile perioden, ville me ha venta ei uniform fordeling av stolar i formrommet. Det observerer me ikkje. Nærmaste-nabo-distansen i det tredimensjonale (høgde, breidde, djupn)-rommet har ein variasjonskoeffisient på 5,4; om lag femten gonger over det ein Poisson-nullhypotese ville predikert. Stolane klumpar seg. Dei samlar seg i visse regionar og let andre stå tomme.

\par\addvspace{2pt}\noindent\centerline{%
  \includegraphics[width=\linewidth]{fig-A.6.1-uniformitet.pdf}}
\par\addvspace{4pt}

Figuren viser gjensidig informasjon mellom kvar geometrisk dimensjon og to kandidatprediktorar: stilperiode og primærmateriale. Stilperiode forklarer 0,48--0,92 bits per dimensjon; materiale forklarer 0,23--0,42 bits. Den kulturelle merkelappen vinn over den fysiske eigenskapen med ein faktor på 1,9 til 2,3 på alle fire dimensjonar samstundes. Materialet er den faktiske mekaniske avgrensninga; stilperioden ber ingen ergonomisk bodskap. Likevel predikerer stilen betre. Dette er den fyrste kvantitative testen av denne påstanden over eit fleirhundreårig korpus.

\section*{Landskapet har fleire toppar}

\par\addvspace{2pt}\noindent\centerline{%
  \includegraphics[width=\linewidth]{fig-A.6.16-fitnesslandskap.pdf}}
\par\addvspace{4pt}

KDE-tettleiken over alle stolar med (breidde, høgde)-data avdekkjer to klare lokale maksimum: éin ved (B \msym{≈} 50, H \msym{≈} 91) cm; den tradisjonelle høgrygga stolen; og éin ved (B \msym{≈} 50, H \msym{≈} 45) cm; den modernistiske låge stolen. Dei tomme dalane mellom toppane er forbodne regionar. Postulatet om éin global optimum kollapsar her; tilpassingslandskapet har minst to stabile attraktorar.

\section*{Det empiriske kraftfeltet}

\par\addvspace{2pt}\noindent\centerline{%
  \includegraphics[width=\linewidth]{fig-A.6.19-kraftfelt.pdf}}
\par\addvspace{4pt}

Tilpassingslandskapet er ikkje ein metafor; det er ei empirisk målbar topologi. Kraftfeltet; den negative gradienten av den logaritmiske tettleiken; er overlagt som eit pilfelt. Kvar piltipp peikar mot den næraste haugen, slik ein agent vil glide om ho følgjer den lokale gradienten. Tre lokale maksimum vert oppdaga automatisk: nyklassisisme, historisme og modernisme. Den turkise bana frå 1500 til 2025 er periodesentroidane; heile formhistoria lagd over landskapet som ei vandring.

\section*{Stilar er gradientar}

\par\addvspace{2pt}\noindent\centerline{%
  \includegraphics[width=\linewidth]{fig-A.6.3-silhouette.pdf}}
\par\addvspace{4pt}

Silhuett-skoren for stilperiode er \msym{−}0,34. Det tyder at gjennomsnittsstolen i ein stilperiode ligg nærmare nabostilane enn dei eigne medkandidatane. 21 av 25 stilar har negativ skore. Stilperiodane er ikkje topologiske klynger; dei er kontinuerlege gradientar i morforommet. Grensene mellom dei svinn; diffusheita er absolutt.

\section*{Kanalisering}

\par\addvspace{2pt}\noindent\centerline{%
  \includegraphics[width=\linewidth]{fig-A.6.2-kanalisering.pdf}}
\par\addvspace{4pt}

Variasjonskoeffisienten over seks meshtrekk strekkjer seg over to storleiksordnar. Sphericity er det mest kanaliserte trekket (CV = 0,074); hylster-volumet det friaste (CV = 9,52). Spreiinga er 128 gonger. Eit funksjonelt avgrensa formrom ville hatt jamn varians. At nokre dimensjonar er låst og andre nær fri, er ein direkte signatur av trykkhierarkiet i tilpassingslandskapet.

\section*{Formrommet ekspanderer}

\par\addvspace{2pt}\noindent\centerline{%
  \includegraphics[width=\linewidth]{fig-A.6.4-ekspansjon.pdf}}
\par\addvspace{4pt}

Det kumulative konvekse hylstervolumet veks monotont gjennom alle tilgjengelege periodar. Frå 79\,392 cm\textsuperscript{3} i fyrste periode til 589\,796 cm\textsuperscript{3}; ein 7 gonger ekspansjon utan eit einaste tilbakesteg. Eit fast funksjonsenvelop ville kravd metning. Den monotone veksten stadfestar at det tilstøytande moglege ekspanderer aktivt.

\section*{Materialstraumen}

\par\addvspace{2pt}\noindent\centerline{%
  \includegraphics[width=\linewidth]{fig-A.6.9-materialstraum.pdf}}
\par\addvspace{4pt}

Materialstraumen frå 1500 til 2025 syner klare seleksjonsbølgjer: eik 1500--1700, nøttetre kring 1675, mahogni 1750--1850, modernismens materiale (stål, plast, kryssfiner, aluminium) etter 1900. Bølgjene er ikkje gradvise overgangar mellom funksjonelt overlegne alternativ; kvar bølgje er ein lokal kohort-event.

\section*{Stiavhengigheit}

\par\addvspace{2pt}\noindent\centerline{%
  \includegraphics[width=\linewidth]{fig-A.6.5-mahogni.pdf}}
\par\addvspace{4pt}

I utvalet av norskproduserte stolar frå 1825--1849 brukar 16 av 16 mahogni. I den føregåande perioden er fraksjonen 0 av 16. Lock-in over éin generasjon. Mahogni og valnøtt er biomekanisk utbytbare for sitjemøbel; dette er ikkje funksjonell optimalisering, men stiavhengig seleksjonskollaps.

\section*{Landskapet endrar seg}

\par\addvspace{2pt}\noindent\centerline{%
  \includegraphics[width=\linewidth]{fig-A.6.6-wasserstein.pdf}}
\par\addvspace{4pt}

Wasserstein-1-distansen mellom suksessive periodar er aldri under terskelen 0,5 cm. Kvar suksessiv periode har ei statistisk distinkt fordeling. Menneskekroppen har ikkje endra seg over 500 år; form endrar seg langt fortare enn funksjon.

\section*{Diskontinuitet}

\par\addvspace{2pt}\noindent\centerline{%
  \includegraphics[width=\linewidth]{fig-A.6.15-rekurrens.pdf}}
\par\addvspace{4pt}

Rekurrensmatrisen over alle 19 periodar viser at det modernistiske brotet etter 1900 er absolutt: ingen periode før 1900 liknar nokon etter. Formhistoria gjentek seg ikkje over dette brotet.

\section*{Latente materialsignaturar}

\par\addvspace{2pt}\noindent\centerline{%
  \includegraphics[width=\linewidth]{fig-A.6.18-materiallatens.pdf}}
\par\addvspace{4pt}

Stål held seg geometrisk nær trestolen i 1900 og 1950 (avstand kring 10 cm), men sprang til 38 cm i 2000. Latensperioden er minst femti år; materialet arvar formspråket til det substratet det erstattar, heilt til eigne affordansar pressar det ut i nye regionar. Materialet ber affordansen lenge før agentane realiserer han.

\section*{Evidens}

Av 6 hovudtestar gonger 14 delmengder (museum, periode, stil) = 84 testkombinasjonar, held samtlege 84. Ingen postulat er falsifiserte i materialet. Funna er robuste på tvers av ulike geometriske representasjonar.

\begin{overgang}
Stolen er ikkje eit val. Ho er eit korpus; eit empirisk laboratorium der proposisjonane let seg teste. At rammeverket held for stolar, er ikkje ein garanti for at det held for bygningar, algoritmar eller organismar. Men det er ein naudsynleg fyrste test. Og den testen er bestått.
\end{overgang}
